\documentclass{article}
\usepackage[a4paper,margin=0.75cm,includefoot]{geometry}
\usepackage[utf8]{inputenc}
\usepackage[T5]{fontenc}
\usepackage{url}
\title{Pretrained Vision Models Across Image Domains\\An Office-Home Evaluation}
\author{}
\date{}
\begin{document}
\maketitle

\section{Objective}
This study evaluates how ImageNet-pretrained vision models perform when training and test images belong to different domains. ResNet18 and ResNet50 are trained on real-world photographs and evaluated on photographs, product images, artistic depictions, and clipart with the same object labels. Two strategies are compared: a frozen backbone with a trainable classification head, and fine-tuning of the final residual block and classification head.

\section{Dataset and evaluation protocol}
Office-Home contains 65 object categories across Art, Clipart, Product, and Real-World. The public Flower Labs mirror is used at a fixed revision. Acquisition verified shard sizes and SHA-256 checksums. Image validation identified 15,588 images, no corrupted images, and 412 byte-identical duplicates. Removing duplicates left 15,176 images.

Data are partitioned within each domain and class into approximately 70\% training, 15\% validation, and 15\% test subsets, using seed 2026. The resulting test sets contain 349 Art, 646 Clipart, 646 Product, and 651 Real-World images. All configurations share the same partitions and class mapping.

Real-World is the source domain. Its training subset alone updates model parameters. Checkpoint selection and early stopping use only source-validation macro-F1. Target-domain training and validation subsets are unused. Final evaluation uses the held-out test subset of each domain. This is a custom source-only domain generalization protocol, not the standard Office-Home domain adaptation benchmark.

\section{Model training}
Both architectures use torchvision's \texttt{IMAGENET1K\_V1} weights and a new 65-class linear head. Linear probing trains only the head. Partial fine-tuning trains \texttt{layer4} and the head while retaining fixed BatchNorm running statistics. No other backbone parameters are updated.

Training images undergo random resized cropping and horizontal flipping. Validation and test images are resized to a shorter edge of 256 pixels and center-cropped. All inputs are $224\times224$ pixels and normalized using ImageNet statistics. Training uses cross-entropy, AdamW, weight decay $10^{-4}$, head learning rate $10^{-3}$, and block learning rate $10^{-4}$. The scheduler uses cosine annealing; gradient norms are clipped to 1.0.

The completed run used a Tesla T4, batch size 32, two DataLoader workers, FP16 mixed precision, and training seed 42. Maximum training lengths are 15 epochs for linear probing and 20 for partial fine-tuning. Early stopping occurs after five epochs without a source-validation macro-F1 improvement. Selected epochs were 8 for ResNet18 linear probe, 4 for ResNet18 fine-tuning, 14 for ResNet50 linear probe, and 12 for ResNet50 fine-tuning. The four runs completed 13, 9, 15, and 17 epochs, respectively.

\section{Test results}
Accuracy and macro-F1 are computed separately for each test domain. Macro-F1 gives equal weight to each of the 65 class-level F1 scores. Accuracy values in the first table are percentages; macro-F1 values use the 0--1 scale.

\begin{center}
\begin{tabular}{llrrrr}
\hline
Model & Strategy & Real-World & Product & Art & Clipart \\
\hline
ResNet18 & Linear probe & 79.42 & 67.34 & 55.59 & 36.22 \\
ResNet18 & Partial fine-tuning & 74.81 & 60.68 & 49.00 & 33.75 \\
ResNet50 & Linear probe & 82.03 & 72.91 & 62.75 & 37.15 \\
ResNet50 & Partial fine-tuning & 82.64 & 71.83 & 61.60 & 40.87 \\
\hline
\end{tabular}

\medskip
\begin{tabular}{llrrrr}
\hline
Model & Strategy & Real-World & Product & Art & Clipart \\
\hline
ResNet18 & Linear probe & 0.7671 & 0.6550 & 0.5182 & 0.3345 \\
ResNet18 & Partial fine-tuning & 0.7230 & 0.5828 & 0.4423 & 0.3160 \\
ResNet50 & Linear probe & 0.7954 & 0.7150 & 0.5898 & 0.3837 \\
ResNet50 & Partial fine-tuning & 0.8109 & 0.7042 & 0.5781 & 0.4372 \\
\hline
\end{tabular}
\end{center}

The same-domain to cross-domain accuracy difference is $\Delta_d=100(A_s-A_d)$, where accuracy values range from 0 to 1. $A_s$ is held-out Real-World accuracy, and $A_d$ is accuracy on domain $d$. The difference is expressed in percentage points. For ResNet50 linear probe, it is 9.12 points for Product, 19.28 for Art, and 44.88 for Clipart.

\section{Interpretation and limitations}
ResNet50 achieves higher accuracy than ResNet18 under each corresponding training strategy in all four domains. Clipart has the lowest accuracy in every configuration. Product images retain higher accuracy than Art and Clipart within each configuration, although domain differences cannot be attributed solely to visual style because image composition and difficulty also differ.

Fine-tuning does not improve all domains. For ResNet18, it reduces accuracy by 4.61 percentage points on Real-World, 6.66 on Product, 6.59 on Art, and 2.48 on Clipart. For ResNet50, it increases accuracy by 0.61 points on Real-World and 3.72 on Clipart, but reduces it by 1.08 on Product and 1.15 on Art. These observations do not identify the cause of the changes or demonstrate statistical significance.

The experiment uses one training seed, one source domain, and one fixed partition. Near-duplicates are not detected by byte-level deduplication, and overlap with ImageNet pretraining data cannot be ruled out. Full-backbone fine-tuning and dedicated adaptation algorithms are not evaluated. No medical, satellite, or industrial dataset is included. The custom protocol prevents direct comparison with standard domain adaptation benchmark scores.

\section{Reproducibility and conclusion}
All four configurations completed in the submitted Colab run. Saved predictions, test membership, accuracy, macro-F1, source-to-target differences, and validation-selected checkpoint epochs were checked for consistency during repository preparation. This check recomputed metrics from prediction files; it did not retrain models or recompute losses from logits.

The results show substantial performance differences across image domains and no uniform benefit from partial fine-tuning. ResNet50 linear probe has the highest accuracy on Art and Product, while ResNet50 partial fine-tuning has the highest accuracy on Real-World and Clipart. Further seeds and source domains are needed to assess the stability of these comparisons.

\section{Member Contribution Table}
All members collaborated throughout the project; the table lists the main area each member focused on.

\begin{center}
\small
\begin{tabular}{p{0.18\textwidth}p{0.13\textwidth}p{0.61\textwidth}}
\hline
Member & Student ID & Main contribution \\
\hline
Nguyễn Tiến Dũng & 23BA14068 & Team leadership, experiment design, model comparison and project integration. \\
Nguyễn Ngọc Hiếu & 23BA14109 & Dataset preparation, data checking, splitting, reproducibility and repository setup. \\
Vũ Minh Châu & 23BA14028 & Image preprocessing, augmentation and data-loading pipeline. \\
Nguyễn Minh Hiếu & 23BA14105 & ResNet18/ResNet50 setup, training and fine-tuning experiments. \\
Lê Đức Anh & 23BA14005 & Evaluation metrics, cross-domain comparison and result visualization. \\
Hoàng Lê Anh Đức & 23BA14057 & Result checking, documentation and final report preparation. \\
\hline
\end{tabular}
\end{center}

\begin{thebibliography}{9}
\bibitem{officehome} H. Venkateswara, J. Eusebio, S. Chakraborty, and S. Panchanathan. \emph{Deep Hashing Network for Unsupervised Domain Adaptation}. CVPR, pp. 5018--5027, 2017.
\bibitem{mirror} Flower Labs. Office-Home dataset mirror. \url{https://huggingface.co/datasets/flwrlabs/office-home}.
\bibitem{torchvision} PyTorch. Torchvision model documentation. \url{https://docs.pytorch.org/vision/stable/models.html}.
\end{thebibliography}
\end{document}
